\section{Translação da espaçonave}
\label{sec_dinAcoplada_translacao_espaconave}

Em um modelo completo de espaçonave, os graus de liberdade de translação do \gls{cdm} aparecem no vetor de coordenadas generalizadas e se acoplam à atitude da espaçonave e ao movimento das juntas do \gls{mr}; assim, deslocamentos internos de massa podem modificar a posição do \gls{cdm} e, em princípio, perturbar o movimento orbital \cite{wie2008space,featherstone2008}.
Neste trabalho, adota-se a simplificação na qual a translação do \gls{cdm} é descrita pelas equações de \gls{hcw} no referencial \gls{lvlh} (Seção~\ref{sec_hillcw}), enquanto a formulação lagrangiana da dinâmica acoplada considera as variáveis de atitude da espaçonave e as juntas do \gls{mr}; desse modo, o vetor de coordenadas generalizadas do sistema é dado por \eqref{eq_q_generalizado_espaconave_mr}.
A cada instante, as equações do movimento orbital relativo fornecem a posição e a velocidade do \gls{cdm} em \gls{lvlh} como entradas para o modelo de atitude e juntas;
por outro lado, os acoplamentos rotacionais entre espaçonave e manipulador são descritos em \eqref{eq_dinamica_geral}.
Essa hipótese limita a avaliação do efeito de deslocamentos internos de massa sobre a órbita, mas preserva a análise do acoplamento atitude--manipulador no modelo dinâmico considerado.
